\documentclass[../../main.tex]{subfiles}% 注意这里的文件路径不能用 ./main.tex ,否则用latexmk编译子文件会报错
\graphicspath{{\subfix{./image/}}{\subfix{../../image/}}} % 指定图片引用路径,后续可以直接使用图片文件名
% 如果图片存放在上级目录,则使用 ../../image/ 作为备用路径

% 例如:
% \begin{figure}[H]
% \centering
% \includegraphics[width=0.3\textwidth]{图.png}
% \caption{}
% \label{figure:图}
% \end{figure}
% 注意:上述\label{}一定要放在\caption{}之后,否则引用图片序号会只会显示??.

\begin{document}

\section{边界对应定理}

\begin{theorem}[边界对应定理]\label{theorem:边界对应定理}
设 \(G\) 是由一条简单闭曲线 \(\Gamma\) 所围成的域, 如果 \(w=f(z)\) 把 \(G\) 双全纯地映为 \(B(0,1)\), 那么 \(f\) 的定义可扩充到 \(\overline{G}\) 上, 使得 \(f \in C(\overline{G})\), 且把 \(\Gamma\) 一一地映为 \(|w|=1\), \(\Gamma\) 关于 \(G\) 的正向对应于 \(f(\Gamma)\) 关于 \(B(0,1)\) 的正向.
\end{theorem}
\begin{proof}
\begin{enumerate}[(1)]
\item 先证明对于任意 \(\zeta \in \partial G, \lim\limits_{\substack{z \to \zeta \\ z \in G}} f(z)\) 存在. 为此只要证明, 如果 \(\lim\limits_{n\to\infty} f(z_n)=a, \lim\limits_{n\to\infty} f(z_n')=b\), 则必有 \(a=b\).
首先, 容易看出, \(a,b\) 都在单位圆周上. 不然的话, 若 \(a \in B(0,1)\), 记 \(w_n = f(z_n)\), 则 \(z_n = f^{-1}(w_n)\). 由于 \(f^{-1}\) 在 \(B(0,1)\) 中连续, 因而有
\begin{align*}
\zeta = \lim\limits_{n\to\infty} z_n = \lim\limits_{n\to\infty} f^{-1}(w_n) = \lim\limits_{w_n \to a} f^{-1}(w_n) = f^{-1}(a).
\end{align*}
这不可能.
现在设 \(a \neq b\), 作分式线性变换 \(T\), 使得 \(T(B(0,1))=B(0,1), T(a)=e^{i\frac{\pi}{4}}, T(b)=e^{5i\frac{\pi}{4}}\). 令 \(g(z)=T(f(z))\), 则 \(g\) 仍然把 \(G\) 双全纯地映为 \(B(0,1)\), 但
\begin{align*}
\lim\limits_{z_n \to \zeta} g(z_n) &= \lim\limits_{z_n \to \zeta} T(f(z_n)) = T(a) = e^{i\frac{\pi}{4}}, \\
\lim\limits_{z_n' \to \zeta} g(z_n') &= \lim\limits_{z_n' \to \zeta} T(f(z_n')) = T(b) = e^{5i\frac{\pi}{4}}.
\end{align*}
对任给的 \(\varepsilon>0\), 必有 \(\delta>0\), 当 \(0<|z_n-\zeta|<\delta, 0<|z_n'-\zeta|<\delta\) 时,
\begin{align*}
|g(z_n)-e^{i\frac{\pi}{4}}|<\varepsilon, \quad |g(z_n')-e^{5\mathrm{i}\frac{\pi}{4}}|<\varepsilon.
\end{align*}
现设 \(z_0 = g^{-1}(0)\), 则 \(z_0 \in G\). 取充分小的 \(\delta>0\), 使得 \(G_\delta = B(\zeta,\delta) \cap G\) 不包含 \(z_0\). 取 \(z_n, z_n' \in G_\delta\), 则 \(g(z_n), g(z_n')\) 分别在以 \(e^{i\frac{\pi}{4}}\) 和 \(e^{5i\frac{\pi}{4}}\) 为中心, 以 \(\varepsilon\) 为半径的圆盘中. 用位于 \(G_\delta\) 中的连续曲线 \(l\) 连接 \(z_n\) 和 \(z_n'\), 则它的像 \(L=g(l)\) 是连接 \(g(z_n)\) 和 \(g(z_n')\) 的 \(B(0,1)\) 中的连续曲线, 它不会经过原点, 因此必和实轴, 虚轴相交, 设交点分别为 \(P\) 和 \(Q\).
用 \(\wideparen{PQ}, \wideparen{P'Q'}, \wideparen{P'Q}\) 分别记 \(PQ\) 相对于实轴, 原点和虚轴对称的弧段, 由这四段弧围成的域记为 \(D\), 显然 \(O \in D\).
令
\begin{align*}
F(w) = (g^{-1}(w)-\zeta)(g^{-1}(\overline{w})-\overline{\zeta}) \cdot (g^{-1}(-w)-\zeta)(g^{-1}(-\overline{w})-\overline{\zeta}),
\end{align*}
则 \(F \in H(B(0,1))\). 我们来估计 \(F\) 在 \(\partial D\) 上的值. 任取 \(w \in \wideparen{PQ}\), 则 \(g^{-1}(w) \in l\), 因此 \(|g^{-1}(w)-\zeta|<\delta\). 如果记 \(M=\sup\limits_{z \in \partial G} |z-\zeta|\), 那么因为 \(g^{-1}(\overline{w}), g^{-1}(-w), g^{-1}(-\overline{w})\) 都属于 \(G\), 因而有
\begin{align*}
|g^{-1}(\overline{w})-\zeta| \leqslant M, \quad
|g^{-1}(-w)-\zeta| &\leqslant M, \quad
|g^{-1}(-\overline{w})-\zeta| \leqslant M.
\end{align*}
于是 \(|F(w)| \leqslant \delta M^3\). 同样道理, 当 \(w\) 属于其他三段弧时也有同样的估计, 即在 \(\partial D\) 上有
\begin{align*}
|F(w)| \leqslant \delta M^3.
\end{align*}
因而由\hyperref[theorem:最大模原理-定理4.5.1]{最大模原理}, 上式在 \(D\) 中也成立. 特别有 \(|F(0)| \leqslant \delta M^3\), 即
\begin{align*}
|g^{-1}(0)-\zeta|^4 \leqslant \delta M^3.
\end{align*}
令 \(\delta \to 0\), 即得 \(g^{-1}(0)=\zeta\), 这不可能. 因而 \(a=b\), 于是 \(\lim\limits_{\substack{z\to\zeta \\ z \in G}} f(z)\) 存在.
\begin{figure}[H]
\centering
\includegraphics[width=0.6\textwidth]{复变图7.1.png}
\caption{}
\label{figure:7_2}
\end{figure}
\item 现在对任意 \(\zeta \in \partial G\), 定义
\begin{align*}
f(\zeta) = \lim\limits_{\substack{z\to\zeta \\ z \in G}} f(z).
\end{align*}
这样, \(f\) 在 \(\overline{G}\) 上都有了定义, 我们证明 \(f \in C(\overline{G})\). 为此, 只需证明对任意 \(\zeta \in \partial G\), \(\lim\limits_{\substack{\xi\to\zeta \\ \xi \in \partial G}} f(\xi) = f(\zeta)\) 就行了. 如\cref{figure:7_3}所示, 固定 \(\zeta \in \partial G\), 对任意 \(\varepsilon>0\), 存在 \(\delta>0\), 当 \(z \in B(\zeta,\delta) \cap G\) 时, \(|f(z)-f(\zeta)|<\varepsilon\). 今取 \(\xi \in B(\zeta,\delta) \cap \partial G\), 对于这个 \(\xi\), 有相应的 \(\delta_1\), 当 \(z \in B(\xi,\delta_1) \cap G\) 时, 有 \(|f(\xi)-f(z)|<\varepsilon\). 今取 \(z \in B(\zeta,\delta) \cap B(\xi,\delta_1) \cap G\), 便有
\begin{align*}
|f(\xi)-f(\zeta)| &\leqslant |f(\xi)-f(z)| + |f(z)-f(\zeta)| < 2\varepsilon.
\end{align*}
这就证明了 \(f \in C(\overline{G})\).
\begin{figure}[H]
\centering
\includegraphics[width=0.3\textwidth]{复变图7.2.png}
\caption{}
\label{figure:7_3}
\end{figure}
\item 前面已经证明, 对每个 \(\zeta \in \partial G, |f(\zeta)|=1\), 这说明 \(f\) 把 \(\partial G\) 映入 \(\partial B(0,1)\). 现在证明, 如果 \(\zeta,\zeta' \in \partial G, \zeta \neq \zeta'\), 那么 \(f(\zeta) \neq f(\zeta')\). 取充分小的 \(\varepsilon>0\), 使得 \(B(\zeta,\varepsilon) \cap B(\zeta',\varepsilon) = \emptyset\), 由于 \(f\) 是单叶的, 因而
\begin{align*}
f(B(\zeta,\varepsilon) \cap G) \cap f(B(\zeta',\varepsilon) \cap G) = \emptyset,
\end{align*}
这就证明了 \(f(\zeta) \neq f(\zeta')\).
由于 \(f\) 是 \(\overline{G}\) 上的一一连续映射, 所以 \(f(\overline{G})\) 是紧集. 注意到 \(B(0,1) \subset f(\overline{G}) \subset \overline{B(0,1)}\), 即知 \(f(\overline{G}) = \overline{B(0,1)}\). 由\cref{proposition:复变例1.7}便知 \(f^{-1}\) 也是 \(\overline{B(0,1)}\) 上的一一连续映射.
\item 最后证明 \(f\) 保持边界的方向不变. 在 \(\partial G\) 上沿着关于 \(G\) 的正方向取三点 \(z_1,z_2,z_3\), 它们在 \(f\) 下的像分别为 \(w_1,w_2,w_3\), 半径 \(Ow_1, Ow_2, Ow_3\) 的原像分别是曲线弧段 \(\wideparen{z_0z_1}, \wideparen{z_0z_2}, \wideparen{z_0z_3}\) (\cref{figure:7_4}). 由于 \(f'(z_0) \neq 0\), \(f\) 在 \(z_0\) 处具有保角性, 而且保持角度的方向不变, 因此 \(w_1, w_2, w_3\) 也是沿着关于 \(B(0,1)\) 的正方向.
\begin{figure}[H]
\centering
\includegraphics[width=0.6\textwidth]{复变图7.3.png}
\caption{}
\label{figure:7_4}
\end{figure}
\end{enumerate}
\end{proof}

\begin{theorem}
设 \(G\) 和 \(D\) 分别是由可求长简单闭曲线 \(\gamma\) 和 \(\Gamma\) 围成的域, 如果 \(f \in H(G) \cap C(\overline{G})\), 且把 \(\gamma\) 一一地映为 \(\Gamma\), 那么 \(w=f(z)\) 把 \(G\) 一一地映为 \(D\), 并且使 \(\gamma\) 关于 \(G\) 的正向对应于 \(\Gamma\) 关于 \(D\) 的正向.
\end{theorem}
\begin{proof}
任取 \(w_0 \in D\), 先证明 \(f(z)-w_0\) 在 \(G\) 中有且只有一个零点. 由辐角原理, \(f(z)-w_0\) 在 \(G\) 中的零点个数为
\begin{align*}
N = \frac{1}{2\pi} \Delta_\gamma \operatorname{Arg}(f(z)-w_0)  = \frac{1}{2\pi} \Delta_\Gamma \operatorname{Arg}(w-w_0) = \pm 1.
\end{align*}
右端等于 1 是因为 \(w_0 \in D\) 之故. 但因 \(N\) 是零点的个数, 不能取 \(-1\), 故有 \(N=1\). 这就证明了 \(f(z)-w_0\) 在 \(G\) 内有且只有一个零点, 且当 \(z\) 沿 \(\gamma\) 的正向转一圈时, \(w\) 也沿 \(\Gamma\) 的正向转一圈. 容易看出, 当 \(w_0 \notin \overline{D}\) 时, \(N=0\), 即 \(f\) 不会把 \(G\) 中的点映到 \(D\) 的外部去. 当 \(w_0 \in \Gamma\) 时, 如果存在 \(z_0 \in G\), 使得 \(f(z_0)=w_0\), 那么 \(w_0\) 必须是 \(f(G)\) 的内点, 这不可能. 这就证明了定理的断言.
\end{proof}

\begin{example}
问 \(w=z^2\) 把由圆周 \(|z-\frac{1}{2}|=\frac{1}{2}\) 围成的域变成什么样的域?
\end{example}
\begin{solution}
从\cref{figure:7_5}可以看出, 圆周 \(|z-\frac{1}{2}|=\frac{1}{2}\) 的极坐标方程为 \(r=\cos\theta\) (\(-\frac{\pi}{2}\leqslant\theta\leqslant\frac{\pi}{2}\)), 这里, \(z=re^{i\theta}\). 若记 \(w=\rho e^{i\varphi}\), 则 \(w=z^2=r^2e^{2i\theta}\), 所以 \(\rho=r^2, \varphi=2\theta\), 故 \(r=\cos\theta\) 经变换后的极坐标方程为 \(\rho=\frac{1}{2}(1+\cos\varphi)\). 这是一条心脏线(图 7.5(b)). 由定理 7.3.2, \(w=z^2\) 把圆周 \(|z-\frac{1}{2}|=\frac{1}{2}\) 的内部一一地映成上述心脏线的内部.
\begin{figure}[H]
\centering
\includegraphics[width=0.8\textwidth]{复变图7.4.png}
\caption{}
\label{figure:7_5}
\end{figure}
\end{solution}



\end{document}